\section{Methods for Markov chains}
\label{m:sec:methods}

Having set out the processes, we turn to the techniques for extracting information
from them. The organising idea is the same throughout: \emph{condition on one step,
and solve the self-consistent equation that results}. In discrete time the step is
a generation and the equation is a functional iteration; in continuous time the
step is the first jump and the equation is a differential one; for generating
functions the step is an infinitesimal interval and the equation is a partial
differential equation.

\subsection{Discrete-time methods}
\label{m:sec:dtmethods}

The techniques below reduce every question about the Galton--Watson process of
\cref{m:sec:gw} to a question about the orbit of the quadratic map
\cref{m:eq:SnIter}. They are set out here because the same three ideas ---
conditioning on one step, generating functions, and functional iteration --- run
through the continuous-time work that follows; the orbit itself is analysed in
\ChA.

\subsubsection{First-step analysis}
\label{m:sec:firststep}

When the question is a hitting probability or a mean absorption time rather than a
full time course, conditioning on the first step replaces the evolution problem by
an algebraic one, because after that step the chain starts afresh from wherever it
landed.

\begin{proposition}[First-step principle]
\label{m:prop:firststep}
Let $\{X_n\}$ be a discrete-time Markov chain with transition matrix $P$ and
absorbing set $\mathcal A$, and let $h_i$ be the probability of ever reaching
$\mathcal A$ from state $i$. Then $h_i=1$ for $i\in\mathcal A$, and for every
transient $i$,
\begin{equation}
\label{m:eq:firststepdiscrete}
h_i = \sum_{j}P_{ij}\,h_j .
\end{equation}
The same conditioning applied to a continuous-time chain with generator $Q$ gives
\begin{equation}
\label{m:eq:firststep}
h_i = \sum_{j\neq i}\frac{q_{ij}}{q_i}\,h_j ,
\end{equation}
and the expected time to absorption satisfies
$k_i = 1/q_i + \sum_{j\neq i}(q_{ij}/q_i)k_j$.
\end{proposition}

\begin{proof}
Absorption occurs eventually if and only if it occurs from the state reached at the
first step; conditioning on that state and using the Markov property gives
\cref{m:eq:firststepdiscrete}. In continuous time the chain waits an
$\mathrm{Exp}(q_i)$ time and then, by \cref{m:prop:race}, jumps to $j$ with
probability $q_{ij}/q_i$ independently of the waiting time, which gives
\cref{m:eq:firststep}; adding the expected holding time $1/q_i$ to the same
decomposition gives the equation for $k_i$.
\end{proof}

The holding times have vanished from \cref{m:eq:firststep}; mean hitting times do
depend on the rates, which is why $1/q_i$ appears in the companion system. The
gambler's-ruin formula \cref{m:eq:gamblersruin} is the two-term instance of
\cref{m:eq:firststepdiscrete}.

For a branching process the first-step decomposition produces something better than
a linear system. Because the offspring of the founder generate independent
subpopulations, conditioning on the first generation gives a \emph{fixed-point
equation}: the ultimate extinction probability $s$ from a single individual
satisfies
\begin{equation}
\label{m:eq:extinctionfixedpoint}
s = \phi(s) ,
\end{equation}
$\phi$ being the offspring generating function.

\subsubsection{Generating functions}
\label{m:sec:pgfmethod}

For a \rv{} $X$ taking values in $\{0,1,2,\dots\}$ with $p_k=\pr{X=k}$, the
\pgf{} is
\begin{equation}
\label{m:eq:pgfdef}
\phi(z) = \ex{z^X} = \sum_{k=0}^{\infty}p_kz^k ,
\end{equation}
convergent at least for $|z|\le1$. The series determines the distribution, since
$p_k=\phi^{(k)}(0)/k!$, and moments are read off at $z=1$:
\begin{equation}
\label{m:eq:pgfmoments}
\phi(1)=1,
\qquad
\phi'(1) = \ex{X},
\qquad
\phi''(1) = \ex{X(X-1)} .
\end{equation}
If $X$ and $Y$ are independent then $\phi_{X+Y}=\phi_X\phi_Y$, so sums become
products; and if $N$ is an independent count with \pgf{} $\psi$, the random sum
$X_1+\dots+X_N$ of independent copies of $X$ has \pgf{} $\psi\circ\phi$, so
\emph{composition} of generating functions corresponds to compounding.

For a pair of counts the definition extends to
\begin{equation}
\label{m:eq:pgfbivariate}
G(x,y) = \ex{x^Xy^Y} = \sum_i\sum_j\pr{X=i,\ Y=j}\,x^iy^j ,
\end{equation}
with marginals recovered by setting the other variable to one. The same
construction with $n$ arguments handles vector-valued states and is the prepared
ground for \fwd{multitype}{the chapter on multi-type processes}; it is not
developed here.

\subsubsection{Functional iteration}
\label{m:sec:iteration}

Applying \cref{m:eq:gwcomposition} repeatedly, the generating function of $Z_n$ is
the $n$-fold composite $\phi_n=\phi^{\circ n}$. For the binary law
\cref{m:eq:pgfbinary} this is the quadratic recursion \cref{m:eq:SnIter} for the
survival probability, obtained by evaluating at $z=0$ and taking complements.

Everything the binary process does is contained in the orbit of that recursion
--- whether $S_n$ reaches zero, how quickly it does so, and what constant stands
in front of the decay when it does --- and following it is the business of \ChA.
Here the recursion is needed only as the denominator of the conditional mean of
\cref{m:sec:condmeans}.

\subsection{Continuous-time methods}
\label{m:sec:ctmethods}

The continuous-time methods are presented on the linear birth--death process of
\cref{m:sec:bd} throughout.

\subsubsection{Forward and backward equations}
\label{m:sec:backward}

Differentiating the Chapman--Kolmogorov identity
$p_{ij}(t+\Delta t)=\sum_kp_{ik}(t)p_{kj}(\Delta t)$ at $\Delta t=0$ gives the
\emph{forward} equations,
\begin{equation}
\label{m:eq:forward}
\ddt{p_{ij}(t)} = \sum_kp_{ik}(t)\,q_{kj},
\qquad\text{that is}\qquad
\dda{P}{t} = PQ ,
\end{equation}
which condition on the last jump before time $t$. Differentiating in the other slot
gives the \emph{backward} equations,
\begin{equation}
\label{m:eq:backwardeq}
\dda{P}{t} = QP ,
\end{equation}
which condition on the first jump after time $0$ --- first-step analysis in
differential form. The two carry the same information but suit different questions:
forward equations track the distribution as it spreads, backward equations track a
functional of the whole future and are the natural setting for absorption
probabilities and for the generating function of a branching process.

Written out for a fixed initial state, with $p_j(t)$ the probability of occupying
state $j$, \cref{m:eq:forward} becomes the \emph{master equation}
\begin{equation}
\label{m:eq:master}
\ddt{p_j(t)} = \sum_{k\neq j}q_{kj}\,p_k(t) - q_j\,p_j(t) :
\end{equation}
probability flows into $j$ from every state that can reach it and out of $j$ at
total rate $q_j$.

For the birth--death rates \cref{m:eq:bdrates} the master equation is
\begin{equation}
\label{m:eq:bdmaster}
\ddt{p_n(t)} = \lambda(n-1)p_{n-1}(t) + \mu(n+1)p_{n+1}(t) - (\lambda+\mu)n\,p_n(t),
\end{equation}
an infinite coupled system, and on its own it is not much use. The step that makes
it tractable is to collapse it into a single equation for a generating function.
Multiply \cref{m:eq:bdmaster} by $z^n$, sum over $n\ge0$, and shift indices so that
every sum is again of the form \cref{m:eq:pgfdef}; the factors of $n$ produced by
the linear rates become derivatives $\partial_z$, and the three terms become
$\lambda z^2\partial_zG$, $\mu\partial_zG$ and $-(\lambda+\mu)z\partial_zG$. Hence
\begin{equation}
\label{m:eq:bdpde}
\pddt{G} = \bigl(\lambda z^2-(\lambda+\mu)z+\mu\bigr)\pdd{G}{z}
= (\lambda z-\mu)(z-1)\pdd{G}{z},
\qquad
G(z,0)=z^N ,
\end{equation}
for a process started from $N$ individuals. This is the template for every
generating-function calculation in this thesis. The quadratic coefficient
factorises with roots $z=1$ and $z=\mu/\lambda$, and the same factorisation
reappears in \cref{m:app:absbirthdeath}.

Solving \cref{m:eq:bdpde} along its characteristics gives the classical result of
Kendall \cite{Kendall1948},
\begin{equation}
\label{m:eq:bdsolution}
G(z,t) = \lrb{\frac{\mu(1-z)-(\mu-\lambda z)\ee^{-(\lambda-\mu)t}}
{\lambda(1-z)-(\mu-\lambda z)\ee^{-(\lambda-\mu)t}}}^{\!N},
\qquad\lambda\neq\mu .
\end{equation}

\subsubsection{Mean population size}
\label{m:sec:mean}

Differentiating \cref{m:eq:bdsolution} at $z=1$ is possible and unpleasant. It is
quicker to take the mean of the master equation directly: multiplying
\cref{m:eq:bdmaster} by $n$ and summing, the terms telescope to
\begin{equation}
\label{m:eq:meanode}
\ddt{}\ex{X_t} = (\lambda-\mu)\ex{X_t},
\qquad
\ex{X_0}=N,
\end{equation}
whence
\begin{equation}
\label{m:eq:bdmean}
\ex{X_t} = N\ee^{(\lambda-\mu)t} .
\end{equation}
The mean obeys a closed equation --- exactly the deterministic rate equation one
would have written down without any probability at all --- and this is a special
feature of \emph{linear} rates. It is also the reason the mean is uninformative
here: \cref{m:eq:bdmean} is the same whether the process is a near-certain
extinction with a rare enormous excursion or a steady moderate growth.

\subsubsection{Variance and higher moments}
\label{m:sec:variance}

The second moment is obtained the same way. Multiplying \cref{m:eq:bdmaster} by
$n^2$ and summing gives
\begin{equation}
\label{m:eq:secondmomentode}
\ddt{}\ex{X_t^2} = 2(\lambda-\mu)\ex{X_t^2} + (\lambda+\mu)\ex{X_t},
\end{equation}
a linear equation driven by the known mean \cref{m:eq:bdmean}. Solving with the
integrating factor $\ee^{-2(\lambda-\mu)t}$ and subtracting $\ex{X_t}^2$ gives, for
$\lambda\neq\mu$,
\begin{equation}
\label{m:eq:bdvariance}
\var{X_t} = N\,\frac{\lambda+\mu}{\lambda-\mu}\,
\ee^{(\lambda-\mu)t}\lrb{\ee^{(\lambda-\mu)t}-1},
\end{equation}
and in the critical case $\lambda=\mu$ the limit is $\var{X_t}=2N\lambda t$. In the
subcritical case $\lambda<\mu$ both mean and
variance decay, but the variance decays more slowly --- the coefficient of
variation grows, and the process becomes relatively \emph{more} random as it dies.
At criticality the mean is constant and the variance grows without bound, which is
the analytic signature of the heavy tail that \ChA{} examines in discrete time.

The pattern generalises. Multiplying the master equation by $n^{\varrho}$ produces
an equation for the $\varrho$th moment driven by the lower ones, so the moments can
be found recursively without ever solving \cref{m:eq:bdpde}. For a process with
nonlinear rates the moment equations do not close, and the generating-function
route of \cref{m:sec:moc} becomes not merely convenient but necessary.

\subsubsection{Hitting probabilities}
\label{m:sec:hitting}

By \cref{m:prop:firststep} the probability $h_i$ of ever reaching a target set from
state $i$ satisfies \cref{m:eq:firststep}, which involves only the ratios
$q_{ij}/q_i$. For the birth--death rates \cref{m:eq:bdrates} those ratios are
independent of the state,
\begin{equation}
\label{m:eq:bdjumpratios}
\frac{q_{n,n+1}}{q_n} = \frac{\lambda n}{(\lambda+\mu)n} = \frac{\lambda}{\lambda+\mu},
\qquad
\frac{q_{n,n-1}}{q_n} = \frac{\mu}{\lambda+\mu},
\end{equation}
so the embedded chain is the simple random walk of \cref{m:sec:randomwalk} with
$q=\lambda/(\lambda+\mu)$, and \cref{m:eq:gamblersruin} answers the question
immediately. With absorbing barriers at $0$ and $M$, the probability of reaching
$M$ before extinction, starting from $i$ individuals, is
\begin{equation}
\label{m:eq:bdruin}
h_i = \frac{1-(\mu/\lambda)^i}{1-(\mu/\lambda)^M}
\qquad(\lambda\neq\mu),
\qquad
h_i = \frac{i}{M}
\qquad(\lambda=\mu),
\end{equation}
since $\theta = (1-q)/q = \mu/\lambda$. Letting $M\to\infty$ gives the probability
of unbounded growth, $1-(\mu/\lambda)^i$ when $\lambda>\mu$ and $0$ otherwise,
which is the complement of the extinction probability found independently in
\cref{m:sec:extinctionprob}. The two routes agreeing is a useful check on both.

Mean hitting times behave differently. They
depend on the rates and not merely on their ratios, and at criticality
$\lambda=\mu$ the mean time to absorption is infinite even though absorption is
certain --- the continuous-time counterpart of the phenomenon derived in full for
the critical Galton--Watson process in \ChA.

\subsubsection{Extinction probabilities}
\label{m:sec:extinctionprob}

Extinction by time $t$ is the event $X_t=0$, so its probability is the generating
function evaluated at the origin: $p_0(t)=G(0,t)$. From \cref{m:eq:bdsolution},
\begin{equation}
\label{m:eq:p0t}
p_0(t) =
\begin{cases}
\displaystyle\lrb{\frac{\mu-\mu\ee^{(\mu-\lambda)t}}{\lambda-\mu\ee^{(\mu-\lambda)t}}}^{\!N},
& \lambda\neq\mu,\\[14pt]
\displaystyle\lrb{\frac{\lambda t}{1+\lambda t}}^{\!N},
& \lambda=\mu,
\end{cases}
\end{equation}
the critical case following by taking $\mu\to\lambda$ in the first line. Letting
$t\to\infty$ recovers the threshold: extinction is certain when $\lambda\le\mu$ and
occurs with probability $(\mu/\lambda)^N$ when $\lambda>\mu$, which is the
complement of \cref{m:eq:bdruin} with $M\to\infty$, as it must be.

\Cref{m:fig:extinction} plots \cref{m:eq:p0t}, alongside the law it leaves behind.

\begin{figure}[ht]
\centering
\includegraphics[width=0.98\textwidth]{extinction_and_law.pdf}
\caption{(a) The extinction probability \cref{m:eq:p0t} against time, with $\mu=1$
and $\lambda=0.6,\,1,\,1.4$; solid curves start from $N=1$, dashed from $N=5$.
Below and at criticality every curve climbs to one, however large the founding
cohort; above it the curves level off at $(\mu/\lambda)^N$, marked dotted for
$N=1$, and a larger cohort buys a permanently smaller extinction probability rather
than merely a delay.
(b) The critical case in detail, from one founder: the law of $X_t$ on the positive
states at $t=1,3,10$, on a logarithmic scale. Absorbing mass accumulates at zero
--- $p_0$ rises from $0.50$ to $0.91$ --- while the surviving tail flattens and
stretches. The mean is exactly $1$ at all three times.}
\label{m:fig:extinction}
\end{figure}

The survival probability is $S^{(N)}(t)=1-p_0(t)$, and for a single founder,
$N=1$, \cref{m:eq:p0t} rearranges to
\begin{equation}
\label{m:eq:ctsS}
S(t) = \frac{\lambda-\mu}{\lambda-\mu\ee^{(\mu-\lambda)t}}
= \frac{\mu-\lambda}{\mu}\,\ee^{-(\mu-\lambda)t}
\lrb{1-\frac{\lambda}{\mu}\ee^{-(\mu-\lambda)t}}^{-1}
\;\sim\;\frac{\mu-\lambda}{\mu}\,\ee^{-(\mu-\lambda)t}
\end{equation}
as $t\to\infty$, in the subcritical case $\mu>\lambda$. The last form is the one
that matters below: the survival probability decays exponentially with a definite
\emph{amplitude} $(\mu-\lambda)/\mu$.

\subsubsection{Conditional means}
\label{m:sec:condmeans}

Absorption is often inevitable --- for $\lambda<\mu$, and for the whole subcritical
branching range --- and yet a process on its way to absorption need not be
featureless. In many cases it settles, long before it is absorbed, into a
distribution that no longer changes with time: a distribution describing the
population conditional on its not yet having been absorbed. The phenomenon is
\emph{quasi-stationarity} and the limit is the Yaglom limit \cite{Yaglom1947};
modern accounts are Méléard and Villemonais \cite{Meleard2012} and Collet,
Martínez and San Martín \cite{Collet2013}.

\begin{definition}[Quasi-stationary and Yaglom laws]
\label{m:def:qsd}
Let $\{X_t\}$ be a Markov process with absorbing state $0$. A probability
distribution $\nu$ on the non-absorbing states is \emph{quasi-stationary} if
$\pr{X_t=j\mid X_t\neq0}=\nu_j$ for all $t\ge0$ and $j\neq0$ whenever
$X_0\sim\nu$. If, from a specified initial state,
$\pr{X_t=j\mid X_t\neq0}\to\nu_j$, then $\nu$ is the corresponding \emph{Yaglom
limit}.
\end{definition}

Existence of the Yaglom limit for the processes used here is classical, from Yaglom
\cite{Yaglom1947} with modern treatments in \cite{Harris1963,Athreya1972,Meleard2012},
and is taken from that literature rather than re-proved. What is computed below is
the limiting conditional mean, explicitly, in both time scales. The distinction
matters: convergence of moments is consistent with the Yaglom theorem but does not
on its own establish convergence of the full conditional law.

\paragraph{The continuous-time constant.}
Since $X_t=0$ contributes nothing to its own expectation,
\begin{equation}
\label{m:eq:condit}
\ex{X_t\mid X_t>0} = \frac{\ex{X_t}}{S^{(N)}(t)} ,
\end{equation}
and for $N=1$ substituting \cref{m:eq:bdmean} and \cref{m:eq:ctsS} gives, exactly,
\begin{equation}
\label{m:eq:ctscondmean}
\ex{X_t\mid X_t>0}
= \frac{\ee^{(\lambda-\mu)t}\lrb{\lambda-\mu\ee^{(\mu-\lambda)t}}}{\lambda-\mu}
= \frac{\lambda\ee^{(\lambda-\mu)t}-\mu}{\lambda-\mu} .
\end{equation}
For $\mu>\lambda$ the exponential decays and the limit is a constant,
\begin{equation}
\label{m:eq:ctsqsmean}
\lim_{t\to\infty}\ex{X_t\mid X_t>0} = \frac{\mu}{\mu-\lambda}
= \frac{1}{1-\lambda/\mu}
=: \frac{1}{\Ac} ,
\end{equation}
the reciprocal of the amplitude identified in \cref{m:eq:ctsS}. To compare this with
the discrete-generation model the two parametrisations must be matched, and
\cref{m:prop:race} is what matches them: the probability that the birth clock wins
is $p=\lambda/(\lambda+\mu)$, so $\lambda/\mu=p/(1-p)$, and
\begin{equation}
\label{m:eq:Acts}
\Ac(p) = \frac{1-2p}{1-p} = \frac{2\varepsilon}{1+\varepsilon},
\qquad
\frac{1}{\Ac(p)} = \frac{1-p}{1-2p},
\qquad
\varepsilon := 1-2p .
\end{equation}
The symbol $\varepsilon$ is used for the distance from criticality throughout this
thesis. \Cref{m:eq:Acts} is elementary and complete: $\Ac$ is a rational function of
$p$, decreasing from $\Ac(0)=1$ to $\Ac(\tfrac12^-)=0$, with
$\Ac(p)\to2\varepsilon$ near criticality.

The reason a closed form was available deserves emphasis, because it is exactly
what fails in discrete time. The survival probability \cref{m:eq:ctsS} is elementary
at \emph{every} time, not merely asymptotically, because \cref{m:eq:bdpde} is a
linear first-order equation with elementary characteristics. The amplitude is a
value of that solution; nothing had to accumulate.

\paragraph{The discrete-time constant.}
For the binary Galton--Watson process the same decomposition gives
$\ex{Z_n\mid Z_n>0} = (2p)^n/S_n$. The mean in the numerator is exactly geometric;
the survival probability in the denominator is only asymptotically so, since the
quadratic term in \cref{m:eq:SnIter} becomes negligible once $S_n$ is small, and
the decay therefore carries a constant in front of it,
\begin{equation}
\label{m:eq:Apdef}
S_n\sim A(p)\,(2p)^n,
\qquad\text{that is}\qquad
A(p) = \lim_{n\to\infty}\frac{S_n}{(2p)^n} ,
\end{equation}
the \emph{amplitude} of the decay. Granting the limit --- which is established in
\ChA, along with everything else that is known about $A(p)$ --- the geometric
factors cancel exactly:
\begin{equation}
\label{m:eq:qsmean}
\ex{Z_n\mid Z_n>0}
\;\xrightarrow[n\to\infty]{}\;
\frac{(2p)^n}{A(p)(2p)^n} = \frac{1}{A(p)} .
\end{equation}
The limiting conditional mean of the discrete-generation process is thus the
reciprocal of its amplitude, precisely as \cref{m:eq:ctsqsmean} is the reciprocal
of the continuous-time amplitude $\Ac$. The same argument carried to the second
moment returns a limiting conditional variance which is again a function of $p$
and $A(p)$ alone; it is stated and derived in \ChA, and is not needed here.

\paragraph{The two constants compared, and the handover.}
Setting the two side by side,
\begin{equation}
\label{m:eq:twoconstants}
\Ac(p) = \frac{1-2p}{1-p}\quad\text{(a closed form)},
\qquad
A(p) = \lim_{n\to\infty}\frac{S_n}{(2p)^n}\quad\text{(a limit, and no more)},
\end{equation}
and \cref{m:fig:dtct} compares them.

\begin{figure}[ht]
\centering
\includegraphics[width=0.48\textwidth]{dtctA.png}
\caption{The continuous-time constant $\Ac(p)=(1-2p)/(1-p)$ (green) against
high-precision numerical values of the discrete-time constant $A(p)$ (red points).
Both vanish at $p=\tfrac12$, but at $p=0$ the continuous curve starts at $1$ and
the discrete one at $\tfrac12$, and the two are not related by a constant factor in
between. The black line is a formula obtained by symbolic regression; it tracks the
points closely but is not exact, and \ChA{} explains why no elementary formula
could be.}
\label{m:fig:dtct}
\end{figure}

\ChAcap{} takes up the constant itself. It shows that the limit \cref{m:eq:Apdef}
exists, with $0<A(p)\le\tfrac12$, so that the quasi-stationary mean is never below
$2$; it represents $A(p)$ as a convergent infinite product; it derives the limiting
conditional variance noted above; it identifies $A(p)$ with a value of the Koenigs
linearising function of the logistic map, and proves that function
hypertranscendental, which is why the search for an elementary formula was never
going to succeed; and it gives the near-critical behaviour together with a
practical scheme for computing $A(p)$ across the whole subcritical range. Two of
those facts are all the present chapter needs: that the limit exists, and that it
is computable to arbitrary accuracy by iterating \cref{m:eq:SnIter}, which is how
the values plotted in \cref{m:fig:dtct} were obtained.
